\documentclass[11pt,a4paper]{article}
\usepackage[utf8]{inputenc}
\usepackage[vietnamese]{babel}
\usepackage{amsmath,amssymb}
\usepackage{geometry}
\geometry{a4paper, total={180mm,260mm}, left=15mm, right=15mm, top=18mm, bottom=18mm}
\usepackage{tikz}
\usetikzlibrary{shapes.geometric, arrows.meta, positioning, calc}
\usepackage{tcolorbox}
\tcbuselibrary{skins,breakable}
\usepackage{xcolor}
\usepackage{multicol}
\usepackage{array}
\usepackage{fontawesome}
\usepackage{fancyhdr}

\pagestyle{fancy}
\fancyhf{}
\renewcommand{\headrulewidth}{0.5pt}
\renewcommand{\footrulewidth}{0.5pt}
\fancyhead[L]{\small\bfseries\color{teal}\faBook\ SỔ TAY KIẾN THỨC TOÁN 3}
\fancyhead[R]{\small\bfseries\color{teal}Lớp Toán Thầy Lĩnh -- loptoanthaylinh.vn}
\fancyfoot[C]{\small\color{gray}-- Sổ tay kiến thức: Bài 49. Luyện tập chung các số đến 10 000 --}

\definecolor{primary}{HTML}{005F73}
\definecolor{secondary}{HTML}{0A9396}
\definecolor{accent}{HTML}{EE9B00}
\definecolor{danger}{HTML}{AE2012}
\definecolor{boxbg}{HTML}{F4F6F8}

\begin{document}

% --- BANNER TIÊU ĐỀ ---
\begin{tcolorbox}[
  enhanced,
  colback=white,
  colframe=primary,
  arc=3mm,
  boxrule=1.2pt,
  left=4mm, right=4mm, top=2.5mm, bottom=2.5mm,
  drop shadow=black!10!white
]
  \centering
  {\Large\bfseries\color{primary}\faBookmark\ BÀI 49: LUYỆN TẬP CHUNG CÁC SỐ ĐẾN 10 000}\\[1.5mm]
  {\small\itshape\color{black!75}Đại tổng kết cấu tạo số có 4 chữ số, so sánh dãy số, chữ số La Mã và kĩ năng làm tròn số trên 1 trang A4}
\end{tcolorbox}

\vspace{-2mm}


\noindent
\begin{minipage}[t]{0.49\textwidth}\vspace{0pt}
\begin{tcolorbox}[
  enhanced,
  colback=blue!4!white,
  colframe=blue!75!black,
  arc=2mm, boxrule=0.8pt,
  left=2mm, right=2mm, top=1.8mm, bottom=1.8mm,
  title={\footnotesize\bfseries\color{white}\faKey\ 1. Ma trận 4 kiến thức cốt lõi}
]
{\footnotesize
\begin{itemize}\setlength{\itemsep}{1.5pt}\setlength{\leftskip}{-2.5mm}
  \item \textbf{Cấu tạo số 4 chữ số:} Hàng nghìn, trăm, chục, đơn vị.
  \item \textbf{So sánh:} Số nào nhiều chữ số hơn thì lớn hơn; cùng số chữ số thì so sánh từ trái sang phải.
  \item \textbf{Chữ số La Mã:} $\text{I}=1, \text{V}=5, \text{X}=10$. Trừ ở bên trái, cộng ở bên phải.
  \item \textbf{Làm tròn số:} Mốc $< 5$ làm tròn xuống, $\ge 5$ làm tròn lên.
\end{itemize}
}
\end{tcolorbox}
\end{minipage}
\hfill
\begin{minipage}[t]{0.49\textwidth}\vspace{0pt}
\begin{tcolorbox}[
  enhanced,
  colback=teal!4!white,
  colframe=teal!75!black,
  arc=2mm, boxrule=0.8pt,
  left=2mm, right=2mm, top=1.8mm, bottom=1.8mm,
  title={\footnotesize\bfseries\color{white}\faSliders\ 2. Kĩ năng lập số từ chữ số cho trước}
]
{\footnotesize
\begin{itemize}\setlength{\itemsep}{1.5pt}\setlength{\leftskip}{-2.5mm}
  \item \textbf{Lập số lớn nhất:} Xếp các chữ số theo thứ tự giảm dần từ hàng nghìn đến hàng đơn vị.
  \item \textbf{Lập số bé nhất:} Chọn chữ số nhỏ nhất khác $0$ ở hàng nghìn, các chữ số còn lại xếp tăng dần.
  \item Đếm số chữ số tròn chục, tròn trăm: hàng đơn vị (và hàng chục) là các chữ số $0$.
\end{itemize}
}
\end{tcolorbox}
\end{minipage}

\vspace{-2.5mm}
% --- KHỐI 3: SƠ ĐỒ TƯ DUY TIKZ ---
\noindent
\begin{tcolorbox}[
  enhanced,
  colback=orange!4!white,
  colframe=orange!80!black,
  arc=2mm, boxrule=0.8pt,
  left=2mm, right=2mm, top=1.5mm, bottom=1.5mm,
  title={\footnotesize\bfseries\color{white}\faSitemap\ SƠ ĐỒ TƯ DUY: HỆ THỐNG KIẾN THỨC CHỦ ĐỀ SỐ ĐẾN 10 000}
]
\centering
\begin{tikzpicture}[
  box/.style={draw, rounded corners=2pt, align=center, font=\scriptsize\bfseries, inner sep=2.5pt, minimum height=0.75cm},
  arr/.style={-{Stealth[length=2mm]}, thick}
]
  \node[box, fill=purple!15, draw=purple!80!black, text width=2.6cm] (q1) at (0, 0) {ĐỌC VIẾT SỐ\\Nghìn $\to$ Đơn vị\\$10\,000$ có 5 chữ số};
  \node[box, fill=teal!15, draw=teal!80!black, text width=2.6cm] (q2) at (3.5, 0) {SO SÁNH THỨ TỰ\\Xét từ trái sang phải\\Tăng dần / Giảm dần};
  \node[box, fill=blue!12, draw=blue!80!black, text width=2.6cm] (q3) at (7.0, 0) {SỐ LA MÃ\\$\text{I}, \text{V}, \text{X}$\\Phạm vi 1 đến 20};
  \node[box, fill=orange!15, draw=orange!80!black, text width=2.6cm] (q4) at (10.5, 0) {LÀM TRÒN SỐ\\Tròn chục $\to$ xét ĐV\\Tròn trăm $\to$ xét Chục};

  \draw[arr, purple!80!black] (q1) -- (q2);
  \draw[arr, teal!80!black] (q2) -- (q3);
  \draw[arr, blue!80!black] (q3) -- (q4);
\end{tikzpicture}
\end{tcolorbox}

\vspace{-2.5mm}
% --- KHỐI 4: VÍ DỤ & MẸO TRÁNH BẪY ---
\noindent
\begin{minipage}[t]{0.49\textwidth}\vspace{0pt}
\begin{tcolorbox}[
  enhanced,
  colback=green!4!white,
  colframe=green!60!black,
  arc=2mm, boxrule=0.8pt,
  left=2mm, right=2mm, top=1.8mm, bottom=1.8mm,
  title={\footnotesize\bfseries\color{white}\faCheckCircle\ Ví dụ thực hành chuẩn mực}
]
{\footnotesize
\textbf{1. Lập số từ bốn chữ số $\{0; 2; 5; 8\}$:}\\
-- Số có 4 chữ số khác nhau \textbf{lớn nhất}: $\mathbf{8520}$.\\
-- Số có 4 chữ số khác nhau \textbf{bé nhất}: $\mathbf{2058}$ (chữ số hàng nghìn phải khác $0$).\\
\textbf{2. Làm tròn phối hợp:}\\
Số học sinh trường Tiểu học là $1465$ bạn.\\
-- Làm tròn đến hàng chục: $\mathbf{1470}$ bạn.\\
-- Làm tròn đến hàng trăm: $\mathbf{1500}$ bạn.
}
\end{tcolorbox}
\end{minipage}
\hfill
\begin{minipage}[t]{0.49\textwidth}\vspace{0pt}
\begin{tcolorbox}[
  enhanced,
  colback=red!4!white,
  colframe=red!70!black,
  arc=2mm, boxrule=0.8pt,
  left=2mm, right=2mm, top=1.8mm, bottom=1.8mm,
  title={\footnotesize\bfseries\color{white}\faExclamationTriangle\ Mẹo tránh bẫy trắc nghiệm}
]
{\footnotesize
\begin{itemize}\setlength{\itemsep}{1.5pt}\setlength{\leftskip}{-2.5mm}
  \item \textbf{Bẫy chọn chữ số 0 đứng đầu hàng nghìn:}\\
  Khi lập số bé nhất, nhiều bạn viết $0258$! Số này chỉ có $3$ chữ số ($258$). Chữ số hàng nghìn bắt buộc phải là chữ số khác $0$ bé nhất!
  \item Không dùng dấu phẩy phân tách hệ nghìn trong văn bản Toán Việt Nam.
\end{itemize}
}
\end{tcolorbox}
\end{minipage}


\end{document}
